\chapter{Complexity — Circuit Eval Machine}
%%% generated by the blueprint pipeline from TCSlib.Complexity.CircuitComplexity.CircuitEvalMachine
\section{Overview}

Declarations of the Lean module \texttt{TCSlib.Complexity.CircuitComplexity.CircuitEvalMachine}.

\section{Declarations}

\begin{definition}[Circuit transition as a machine action]
\lean{BoolCircuit.CircuitEval.LA.toAct}
\label{BoolCircuit.CircuitEval.LA.toAct}
\leanok
\uses{BoolCircuit.CircuitEval.LA, BoolCircuit.CircuitEval.St, Turing.Action}
Given an input-head movement $m\in\{-1,0,1\}$ and a transition record with value-tape
write instruction $v$, mark-tape write instruction $w$, shared work-head movement
$h\in\{-1,0,1\}$, optional output symbol $o$, and optional successor state $q$, form the
action of a two-work-tape Boolean Turing machine that moves its input head by $m$,
applies $v$ to the first work tape and $w$ to the second, moves both work heads by $h$,
emits $o$, and passes to $q$, halting if $q$ is absent.
\end{definition}

\begin{definition}[Pairwise Boolean symbol encoding]
\lean{BoolCircuit.CircuitEval.pairSym}
\label{BoolCircuit.CircuitEval.pairSym}
\leanok
For Booleans $b$ and $b'$, return the nested optional symbol containing $b$ if they are
equal; if they differ and $b$ is false, return the encoding of a blank; otherwise return
no encoding.
\end{definition}

\begin{definition}[Circuit evaluation transition function]
\lean{BoolCircuit.CircuitEval.tr}
\label{BoolCircuit.CircuitEval.tr}
\leanok
\uses{BoolCircuit.CircuitEval.LA.go, BoolCircuit.CircuitEval.LA.rej, BoolCircuit.CircuitEval.LA.toAct, BoolCircuit.CircuitEval.St, BoolCircuit.CircuitEval.lact, BoolCircuit.CircuitEval.pairSym, Turing.Action, Turing.FinTM.controlAction}
For a Boolean flag $e$, define the transition function of a two-work-tape machine on
each state, input symbol, and pair of scanned work-tape symbols as follows. It reads
input symbols in pairs and rejects if a required input symbol is blank or the pair
cannot be encoded by the pairwise Boolean symbol encoding; otherwise it parses the
circuit description, checks gate arguments and gate arities, evaluates each gate, and
eventually emits the computed Boolean. During parsing, malformed encodings are rejected;
while copying, if $e$ is true, any additional input symbol after the copied description
causes rejection, whereas if $e$ is false the machine may stop copying at that point.
\end{definition}

\begin{definition}[Circuit evaluation machine]
\lean{BoolCircuit.CircuitEval.evalTM}
\label{BoolCircuit.CircuitEval.evalTM}
\leanok
\uses{BoolCircuit.CircuitEval.St, BoolCircuit.CircuitEval.tr, Turing.FinTM}
For a Boolean flag $e$, define a finite Turing machine with two work tapes, Boolean tape
symbols, and the circuit-evaluation transition function for $e$. Its initial state
begins the input-reading and circuit-parsing process.
\end{definition}

\begin{definition}[Two-tape projection]
\lean{BoolCircuit.CircuitEval.Internal.tp}
\label{BoolCircuit.CircuitEval.Internal.tp}
\leanok
For functions $v,m:\mathbb{Z}\to\mathrm{Option}(\mathrm{Bool})$, define a function on
two tape indices by returning $v$ at index $0$ and $m$ at index $1$, at every integer
position.
\end{definition}

\begin{definition}[Boolean tape configuration constructor]
\lean{BoolCircuit.CircuitEval.Internal.cf}
\label{BoolCircuit.CircuitEval.Internal.cf}
\leanok
\uses{BoolCircuit.CircuitEval.Internal.tp, BoolCircuit.CircuitEval.St, Turing.Cfg}
Given a finite Boolean input string $z$, an optional machine state $q$, an input-head
position $p$ among the $|z|+2$ positions, two tape contents
$v_t,m_t:\mathbb{Z}\to\mathrm{Option}(\\text{Bool})$, a work-tape head position
$j\in\mathbb{Z}$, and a finite Boolean output string, this construction forms a
configuration of a two-work-tape Turing machine with input $z$, state $q$, input-head
position $p$, both work-tape head positions equal to $j$, and the given output. The
work-tape contents are obtained from $v_t$ and $m_t$ by the construction used to combine
them into the two tapes.
\end{definition}

\begin{definition}[Optional tape-cell write]
\lean{BoolCircuit.CircuitEval.Internal.wr}
\label{BoolCircuit.CircuitEval.Internal.wr}
\leanok
Let $t$ be a tape assigning each integer position either a Boolean value or no value,
and let $j$ be an integer position. Given an optional write instruction whose content is
itself either a Boolean value or no value, the resulting tape is $t$ if there is no
instruction; otherwise it agrees with $t$ away from $j$ and has the instructed content
at $j$.
\end{definition}

\begin{definition}[Input head position]
\lean{BoolCircuit.CircuitEval.Internal.ip}
\label{BoolCircuit.CircuitEval.Internal.ip}
\leanok
For a Boolean string $z$ and an index $i\in\mathbb{N}$ with $i\le |z|$, define the
corresponding input-head position to be $i+1$ among the $|z|+2$ positions.
\end{definition}

\begin{definition}[Marked-cell tape]
\lean{BoolCircuit.CircuitEval.Internal.tapeM}
\label{BoolCircuit.CircuitEval.Internal.tapeM}
\leanok
For a finite list $M$ of natural numbers, define a tape on the integer positions that
contains true at position $c$ exactly when $c$ is nonnegative and its natural-number
value belongs to $M$; every other position is blank.
\end{definition}

\begin{theorem}[One-step configuration update]
\lean{BoolCircuit.CircuitEval.step_cf}
\label{BoolCircuit.CircuitEval.step_cf}
\leanok
\uses{BoolCircuit.CircuitEval.Internal.cf, BoolCircuit.CircuitEval.Internal.tp, BoolCircuit.CircuitEval.Internal.wr, BoolCircuit.CircuitEval.LA, BoolCircuit.CircuitEval.LA.toAct, BoolCircuit.CircuitEval.St, BoolCircuit.CircuitEval.evalTM, BoolCircuit.CircuitEval.tr, Turing.Action.apply, Turing.Cfg.inputSymbol, Turing.MultiTapeTM.step, Turing.moveInputPos}
\difficulty{4}
Let $z$ be a finite bit string, and consider a configuration with state $q$, input-head
position $p$, two tape contents $v$ and $m$, both work heads at the integer position
$j$, and output string $o$. For any input-head movement $i$ and local action $a$, if the
machine transition on this configuration's scanned input symbol and the symbols on its
two tapes equals the action encoded by $a$ with input movement $i$, then one machine
step yields state and input-head position prescribed by $a$ and the clamped movement of
$p$ by $i$; it updates the two tape contents at $j$ according to their respective
writes, moves both work heads by the movement prescribed by $a$, and appends its
optional output symbol to $o$.
\end{theorem}

\begin{theorem}[Two-step circuit-machine transition]
\lean{BoolCircuit.CircuitEval.lstep}
\label{BoolCircuit.CircuitEval.lstep}
\leanok
\uses{BoolCircuit.CircuitEval.Internal.cf, BoolCircuit.CircuitEval.Internal.ip, BoolCircuit.CircuitEval.Internal.tp, BoolCircuit.CircuitEval.Internal.wr, BoolCircuit.CircuitEval.LA.go, BoolCircuit.CircuitEval.Ph, BoolCircuit.CircuitEval.evalTM, BoolCircuit.CircuitEval.lact, BoolCircuit.CircuitEval.pairSym, BoolCircuit.CircuitEval.step_cf, BoolCircuit.CircuitEval.tr, Turing.FinTM.inputSymbol_at, Turing.MultiTapeTM.runFrom, Turing.MultiTapeTM.step, Turing.moveInputPos_pos_of_ne_right}
\difficulty{4}
Let $z$ be a Boolean word, let $i+2\le |z|$, and suppose its entries at positions $i$
and $i+1$ are $b$ and $b'$. If the symbol pair determined by $b,b'$ is $s$, then, for
either setting of the exactness parameter, the circuit-evaluation machine started in
phase $\mathrm{rd1}(\phi)$ at input position $i+1$, with value and mark tapes $v_t,m_t$,
both work heads at $j$, and empty output, reaches after exactly two steps the
configuration at input position $i+3$ whose state, tape contents, head positions, and
output are given by the local action for phase $\phi$, symbol $s$, and scanned tape
values $v_t(j),m_t(j)$.
\end{theorem}

\begin{theorem}[Erasing a finite tape segment]
\lean{BoolCircuit.CircuitEval.rew}
\label{BoolCircuit.CircuitEval.rew}
\leanok
\uses{BoolCircuit.CircuitEval.Internal.cf, BoolCircuit.CircuitEval.Internal.tp, BoolCircuit.CircuitEval.Internal.wr, BoolCircuit.CircuitEval.LA.toAct, BoolCircuit.CircuitEval.St, BoolCircuit.CircuitEval.evalTM, BoolCircuit.CircuitEval.step_cf, BoolCircuit.CircuitEval.tr, Turing.MultiTapeTM.runFrom, Turing.MultiTapeTM.runFrom_succ_eq_step, Turing.MultiTapeTM.runFrom_zero, Turing.moveInputPos_zero}
\difficulty{5}
Let $z$ be a Boolean word, let $q$ and $d$ be machine states, and choose one of the two
work tapes, indexed by $k$. Suppose that for every input symbol and every pair of
scanned work-tape symbols, the transition in state $q$ leaves the input head, output,
and tapes unchanged, moves both work heads left when tape $k$ is occupied, and moves
both right and enters state $d$ when tape $k$ is blank. If tape $k$ is occupied at every
integer cell $c$ with $0\le c<h_0$, and blank at cell $-1$, then starting with both work
heads at $h_0-1$ in state $q$, after $h_0+1$ steps the machine is in state $d$ with both
work heads at zero; the input position, tapes, and output are unchanged. This holds for
every natural number $h_0$ and either setting of the exactness parameter.
\end{theorem}

\begin{definition}[Erasing an interval of tape cells]
\lean{BoolCircuit.CircuitEval.Internal.erase}
\label{BoolCircuit.CircuitEval.Internal.erase}
\leanok
Given a tape $t:\mathbb{Z}\to\mathrm{Option}(\mathrm{Bool})$ and natural numbers $h$ and
$len$, define a new tape that is blank at every integer cell $c$ satisfying $h\le c\le
len$, and agrees with $t$ at all other cells.
\end{definition}

\begin{theorem}[Appending a value to the computation tape]
\lean{BoolCircuit.CircuitEval.app_walk}
\label{BoolCircuit.CircuitEval.app_walk}
\leanok
\uses{BoolCircuit.CircuitEval.Internal.cf, BoolCircuit.CircuitEval.Internal.erase, BoolCircuit.CircuitEval.Internal.tp, BoolCircuit.CircuitEval.Internal.wr, BoolCircuit.CircuitEval.evalTM, BoolCircuit.CircuitEval.step_cf, BoolCircuit.CircuitEval.tr, Turing.FinTM.bufferTape, Turing.FinTM.bufferTape_append, Turing.FinTM.bufferTape_nat, Turing.MultiTapeTM.runFrom, Turing.MultiTapeTM.runFrom_succ_eq_step, Turing.MultiTapeTM.runFrom_zero, Turing.moveInputPos_zero}
\difficulty{5}
Let $e$ be the exactness flag, let $z$ be a Boolean input word, let $v$ be a Boolean
value, and let $p$ be a valid input-head position for $z$. For any lists of previously
computed values and output symbols, and any natural numbers $r,h$ satisfying
$h+r=|\mathrm{vals}|$, starting in the append state with value $v$, input-head position
$p$, value tape storing $\mathrm{vals}$, mark tape $m$, work-head position $h$, and
output $\mathrm{out}$, after $r+1$ machine steps the state is the value-rewind state
returning to the gate-reading phase, the value tape stores $\mathrm{vals}$ followed by
$v$, the mark tape is erased on the integer interval from $h$ through $|\mathrm{vals}|$
inclusive, the input-head position and output are unchanged, and both work-head
positions equal $|\mathrm{vals}|-1$.
\end{theorem}

\begin{definition}[Main phase predicate]
\lean{BoolCircuit.CircuitEval.Internal.main}
\label{BoolCircuit.CircuitEval.Internal.main}
\leanok
\uses{BoolCircuit.CircuitEval.Ph}
For the phases of the circuit-evaluation machine, define a Boolean predicate that is
false on the phases $n0$ and $sep$ and true on every other phase.
\end{definition}

\begin{definition}[Walk phase predicate]
\lean{BoolCircuit.CircuitEval.Internal.walk}
\label{BoolCircuit.CircuitEval.Internal.walk}
\leanok
\uses{BoolCircuit.CircuitEval.Ph}
For the phases of the circuit-evaluation machine, define a Boolean predicate that is
true on phases $u(k,c,a)$, $o$, and $fin(v)$, and false on every other phase.
\end{definition}

\begin{definition}[Evaluation-state invariant]
\lean{BoolCircuit.CircuitEval.Internal.Inv}
\label{BoolCircuit.CircuitEval.Internal.Inv}
\leanok
\uses{BoolCircuit.CircuitEval.AS, BoolCircuit.CircuitEval.Internal.main, BoolCircuit.CircuitEval.Internal.walk}
Let an evaluation state consist of a phase $\phi$, a Boolean list $vals$, a list $M$ of
natural numbers marking tape cells, and a natural-number work-head position $j$. The
invariant requires $j\le |vals|$, every $a\in M$ to satisfy $a<|vals|$, and the
main-phase predicate of $\phi$ to be true; it also requires $j=0$ whenever the
walk-phase predicate of $\phi$ is false.
\end{definition}

\begin{definition}[Admissible transition labels]
\lean{BoolCircuit.CircuitEval.Internal.LAok}
\label{BoolCircuit.CircuitEval.Internal.LAok}
\leanok
\uses{BoolCircuit.CircuitEval.Internal.main, BoolCircuit.CircuitEval.Internal.walk, BoolCircuit.CircuitEval.LA}
Let $w$ be a Boolean, let $v_c$ be an optional Boolean, and let $\ell$ be a transition
label with optional value-tape and mark-tape writes, a work-head move, an optional
emitted Boolean, and an optional next state. The label is admissible precisely when it
makes no value-tape write and either makes no mark-tape write or writes a marked true
value while $v_c$ is present. If there is no next state, an emitted Boolean must be
present; if the next state is a first-read state carrying phase $\phi'$, no Boolean is
emitted, $\mathrm{main}(\phi')$ holds, the head is stationary or moves positively with
$v_c$ present, and if $\mathrm{walk}(\phi')$ is false then $w$ is false and the head is
stationary. If the next state is a value-rewind state carrying a phase $\phi'$, no
Boolean is emitted, $\mathrm{main}(\phi')$ holds, $\mathrm{walk}(\phi')$ is false, the
head moves negatively, and $v_c$ is present; if the next state is an append state, no
Boolean is emitted, the head is stationary, no mark-tape write is made, and $w$ is
false. Every other next state is inadmissible.
\end{definition}

\begin{theorem}[Validity conditions for a phase action]
\lean{BoolCircuit.CircuitEval.lact_ok}
\label{BoolCircuit.CircuitEval.lact_ok}
\leanok
\uses{BoolCircuit.CircuitEval.Internal.LAok, BoolCircuit.CircuitEval.Internal.main, BoolCircuit.CircuitEval.Internal.walk, BoolCircuit.CircuitEval.LA.go, BoolCircuit.CircuitEval.LA.rej, BoolCircuit.CircuitEval.Ph, BoolCircuit.CircuitEval.lact}
\difficulty{5}
Let $\phi$ be a phase with $\text{main}(\phi)=\text{true}$, and let $s,v,m$ be optional
Boolean symbols. The action prescribed by the phase and these symbols writes nothing on
the value tape; its mark-tape write is either absent or writes a mark when $v$ is
present. If it halts, it emits a symbol. If its next state is a read state for phase
$\phi'$, it emits nothing, has $\text{main}(\phi')=\text{true}$, and either does not
move the work heads or moves them right when $v$ is present; moreover, whenever
$\text{walk}(\phi')$ is false, the current phase is not a walking phase and the work
heads do not move. If its next state is a value-tape rewind state for a phase $\phi'$,
it emits nothing, has $\text{main}(\phi')=\text{true}$ and
$\text{walk}(\phi')=\text{false}$, moves the work heads left, and has $v$ present. If
its next state is an append state, it emits nothing, does not move the work heads or
write on the mark tape, and the current phase is not a walking phase.
\end{theorem}

\begin{definition}[Initial configuration for a phase]
\lean{BoolCircuit.CircuitEval.Internal.cfA}
\label{BoolCircuit.CircuitEval.Internal.cfA}
\leanok
\uses{BoolCircuit.CircuitEval.AS, BoolCircuit.CircuitEval.Internal.cf, BoolCircuit.CircuitEval.Internal.ip, BoolCircuit.CircuitEval.Internal.tapeM, BoolCircuit.CircuitEval.St, Turing.Cfg, Turing.FinTM.bufferTape}
Let $z$ be a Boolean input list, let $s$ specify a phase together with a list of
previously computed vertex values, a list of marked cells, and a work-head position $j$,
and let $i$ be a natural number with $i\le |z|$. The associated two-work-tape machine
configuration has input $z$, first-read state carrying the phase of $s$, input-head
position determined by $i$, value tape initialized from the recorded vertex values, mark
tape initialized from the recorded marked cells, work-head position $j$, and empty
output.
\end{definition}

\begin{theorem}[Natural-index tape marks]
\lean{BoolCircuit.CircuitEval.tapeM_nat}
\label{BoolCircuit.CircuitEval.tapeM_nat}
\leanok
\uses{BoolCircuit.CircuitEval.Internal.tapeM, BoolCircuit.CircuitEval.markCell}
\difficulty{1}
For every finite list $m$ of natural numbers and every natural number $j$, the mark tape
whose marked cells are the nonnegative integer positions corresponding to members of
$m$, evaluated at position $j$, equals the mark-cell indicator of whether $j$ belongs to
$m$.
\end{theorem}

\begin{theorem}[Adding a marked tape cell]
\lean{BoolCircuit.CircuitEval.wr_tapeM}
\label{BoolCircuit.CircuitEval.wr_tapeM}
\leanok
\uses{BoolCircuit.CircuitEval.Internal.tapeM, BoolCircuit.CircuitEval.Internal.wr}
\difficulty{4}
For every finite list $m$ of natural numbers and every natural number $j$, writing true
in cell $j$ of the tape whose marked cells are the nonnegative positions corresponding
to members of $m$ produces the tape whose marked positions are exactly the members of
the list obtained by placing $j$ at the beginning of $m$.
\end{theorem}

\begin{theorem}[Bounded simulation of one abstract update]
\lean{BoolCircuit.CircuitEval.macro_step}
\label{BoolCircuit.CircuitEval.macro_step}
\leanok
\uses{BoolCircuit.CircuitEval.AS, BoolCircuit.CircuitEval.Internal.Inv, BoolCircuit.CircuitEval.Internal.cf, BoolCircuit.CircuitEval.Internal.cfA, BoolCircuit.CircuitEval.Internal.erase, BoolCircuit.CircuitEval.Internal.ip, BoolCircuit.CircuitEval.Internal.tapeM, BoolCircuit.CircuitEval.Internal.tp, BoolCircuit.CircuitEval.Internal.wr, BoolCircuit.CircuitEval.absOf, BoolCircuit.CircuitEval.absStep, BoolCircuit.CircuitEval.app_walk, BoolCircuit.CircuitEval.evalTM, BoolCircuit.CircuitEval.lact, BoolCircuit.CircuitEval.lact_ok, BoolCircuit.CircuitEval.lstep, BoolCircuit.CircuitEval.markCell, BoolCircuit.CircuitEval.pairSym, BoolCircuit.CircuitEval.rew, BoolCircuit.CircuitEval.tapeM_nat, BoolCircuit.CircuitEval.wr_tapeM, Turing.FinTM.bufferTape, Turing.FinTM.bufferTape_nat, Turing.MultiTapeTM.runFrom, Turing.MultiTapeTM.runFrom_add}
\difficulty{6}
Let $z$ be a Boolean word, let $s$ satisfy the invariant that its work-head index is at
most the length of its computed-value list, every marked index is smaller than that
length, its phase is a main phase, and a phase outside the walking phases has work-head
index zero. Suppose $i+2\le |z|$, the entries of $z$ at $i,i+1$ are $b,b'$, and their
pair symbol is $\sigma$. Apply the abstract one-symbol update to $s$ using $\sigma$. If
it yields a new abstract state $s'$, then $s'$ satisfies the invariant, its value list
grows by at most one, and the circuit-evaluation machine reaches the configuration
encoding $s'$ at input position $i+3$ within $2|s.\mathrm{vals}|+4$ steps; if it yields
a Boolean value $v$, then within that many steps the machine halts with output $[v]$.
\end{theorem}

\begin{theorem}[Outcome on an absent symbol]
\lean{BoolCircuit.CircuitEval.absStep_none}
\label{BoolCircuit.CircuitEval.absStep_none}
\leanok
\uses{BoolCircuit.CircuitEval.AS, BoolCircuit.CircuitEval.LA.rej, BoolCircuit.CircuitEval.absOf, BoolCircuit.CircuitEval.absStep, BoolCircuit.CircuitEval.lact}
\difficulty{4}
For every abstract evaluation state $s$, there exists a Boolean value $v$ such that
processing the absent symbol yields the outcome $v$, rather than another state.
\end{theorem}

\begin{theorem}[Halting simulation on a delimited word]
\lean{BoolCircuit.CircuitEval.main_run}
\label{BoolCircuit.CircuitEval.main_run}
\leanok
\uses{BoolCircuit.CircuitEval.AS, BoolCircuit.CircuitEval.Internal.Inv, BoolCircuit.CircuitEval.Internal.cfA, BoolCircuit.CircuitEval.absRun, BoolCircuit.CircuitEval.absStep, BoolCircuit.CircuitEval.absStep_none, BoolCircuit.CircuitEval.evalTM, BoolCircuit.CircuitEval.macro_step, BoolCircuit.CircuitEval.pairSym, Turing.MultiTapeTM.runFrom, Turing.MultiTapeTM.runFrom_add, Turing.dbl}
\difficulty{5}
Let $z,x$ be Boolean words, and let $l$ be a Boolean word. Suppose
$z=\mathrm{pre}\,\mathrm{dbl}(l)\,(0,1)\,x$, where $\mathrm{pre}$ is a Boolean word, and
let $s$ satisfy the invariant that its work-head index is at most the length of its
computed-value list, every marked index is smaller than that length, its phase is a main
phase, and a phase outside the walking phases has work-head index zero. For any
$B\in\mathbb N$ with $|s.\mathrm{vals}|+|l|\le B$, the circuit-evaluation machine,
started in the configuration encoding $s$ at the input position immediately after
$\mathrm{pre}$, halts with output equal to the Boolean value obtained by recursively
processing $l$ from $s$, within $(|l|+1)(2B+4)$ steps. This holds for either setting of
the exactness parameter.
\end{theorem}
